% Self-contained report. Compile with pdfLaTeX in Overleaf (twice for contents).
\documentclass[12pt,a4paper]{article}
\usepackage[T1]{fontenc}
\usepackage[utf8]{inputenc}
\usepackage[english]{babel}
\usepackage[margin=2.5cm]{geometry}
\usepackage{xcolor}
\usepackage{listings}
\usepackage{booktabs}
\usepackage{microtype}
\usepackage[hidelinks]{hyperref}
\definecolor{listinggray}{RGB}{247,247,247}
\lstset{basicstyle=\ttfamily\footnotesize,backgroundcolor=\color{listinggray},
 frame=single,rulecolor=\color{gray},breaklines=true,breakatwhitespace=false,
 columns=fullflexible,keepspaces=true,showstringspaces=false,tabsize=4,
 captionpos=b,aboveskip=10pt,belowskip=10pt}
\setlength{\emergencystretch}{3em}
\begin{document}
\begin{titlepage}
\centering
\textbf{Ministry of Education, Culture and Research of the Republic of Moldova}\par\medskip
\textbf{Technical University of Moldova}\par\medskip
\textbf{Faculty of Computers, Informatics and Microelectronics}\par\medskip
\textbf{Department of Software Engineering and Automatic}\par
\vspace{2.5cm}
\textbf{ANDREI ELCISCEV}\par
\textbf{Group FAF-243}\par
\vspace{2cm}
{\Large\textbf{LABORATORY WORK No. 18}\par}
\medskip
{\Large\textbf{BUILD A USER-LEVEL THREAD SWITCH}\par}
\vspace{2cm}
\begin{flushright}
\textbf{Checked by:}\\
Rei\c{t}man Patricia
\end{flushright}
\vfill
\textbf{Chi\c{s}in\u{a}u, 2026}
\end{titlepage}
\pagenumbering{arabic}
\tableofcontents
\newpage
\section{Introduction}
This laboratory work implements cooperative user-level threads in xv6. Each
thread has its own stack and a saved register context. A user-space scheduler
switches between runnable threads by calling an assembly routine that saves one
context and restores another.

The working tree is \texttt{/home/user/os-lab18/xv6-riscv}, copied from the
completed Lab 13 tree. The handout's \texttt{os2026-setup.sh} was absent under
\texttt{/home/user}. The installed RISC-V compiler and QEMU were checked directly,
and the baseline built and booted before the Lab 18 edits.

\section{Environment and Baseline}

\begin{lstlisting}[caption={Baseline and installed toolchain}]
Baseline source: /home/user/os-lab13/xv6-riscv
os2026-setup.sh was not found under /home/user. Direct build and boot checks follow.
$ uname -sr
Linux 6.14.0-37-generic
exit code: 0
$ riscv64-linux-gnu-gcc --version
riscv64-linux-gnu-gcc (Ubuntu 13.3.0-6ubuntu2~24.04.1) 13.3.0
Copyright (C) 2023 Free Software Foundation, Inc.
This is free software; see the source for copying conditions.  There is NO
warranty; not even for MERCHANTABILITY or FITNESS FOR A PARTICULAR PURPOSE.

exit code: 0
$ qemu-system-riscv64 --version
QEMU emulator version 8.2.2 (Debian 1:8.2.2+ds-0ubuntu1.18)
Copyright (c) 2003-2023 Fabrice Bellard and the QEMU Project developers
exit code: 0
\end{lstlisting}

\begin{lstlisting}[caption={Baseline QEMU boot and shell check}]
xv6 kernel is booting

hart 1 starting
hart 2 starting
init: starting sh
$ echo baseline boot OK
baseline boot OK
$ QEMU: Terminated
\end{lstlisting}

\section{Part 1: Implement the Switch}
\subsection{The Saved Context}
The context consists of 14 fields of type \texttt{uint64}. The order matches the
offsets used in the assembly routine. The return address occupies offset 0 and
the stack pointer offset 8. The twelve saved registers \texttt{s0} through
\texttt{s11} occupy offsets 16 through 104. The complete context is 112 bytes.
Each thread has an 8192-byte stack aligned to 16 bytes.

\subsection{Thread Creation}
\texttt{thread\_create} selects a free slot, marks it runnable, puts the function
address in \texttt{ctx.ra}, and sets \texttt{ctx.sp} to the top of the stack.
A small boundary check reports an error if all four slots are occupied.

\subsection{Scheduling and Switching}
The scheduler scans the four slots in round-robin order. When it finds a different
runnable thread, it remembers the current thread, marks the selected thread
running, updates \texttt{current\_thread}, and switches the saved contexts.
Yielding marks the current thread runnable before entering the scheduler.

The main thread uses slot 0, then marks itself free after creating the three
workers. Workers A, B, and C each wait until the other two have started, increment
their own counter 100 times with a yield after each iteration, then mark their
slot free and schedule again.

\subsection{Makefile Changes}
The program is added to \texttt{UPROGS}, and the explicit link rule includes both
its C object and its assembly switch object.

\begin{lstlisting}[caption={New UPROGS entry and explicit link rule}]
	$U/_uthread\

$U/_uthread: $U/uthread.o $U/uthread_switch.o $(ULIB) $U/user.ld
	$(LD) $(LDFLAGS) -T $U/user.ld -o $U/_uthread $U/uthread.o $U/uthread_switch.o $(ULIB)
	$(OBJDUMP) -S $U/_uthread > $U/uthread.asm
\end{lstlisting}

\subsection{Assembly Inspection}
The finished tree built successfully. The following actual object disassembly
shows the save/load pairs generated from \texttt{user/uthread\_switch.S}.

\begin{lstlisting}[caption={Disassembly of the context-switch routine}]
$ riscv64-linux-gnu-objdump -d user/uthread_switch.o

user/uthread_switch.o:     file format elf64-littleriscv


Disassembly of section .text:

0000000000000000 <thread_switch>:
   0:	00153023          	sd	ra,0(a0)
   4:	00253423          	sd	sp,8(a0)
   8:	e900                	sd	s0,16(a0)
   a:	ed04                	sd	s1,24(a0)
   c:	03253023          	sd	s2,32(a0)
  10:	03353423          	sd	s3,40(a0)
  14:	03453823          	sd	s4,48(a0)
  18:	03553c23          	sd	s5,56(a0)
  1c:	05653023          	sd	s6,64(a0)
  20:	05753423          	sd	s7,72(a0)
  24:	05853823          	sd	s8,80(a0)
  28:	05953c23          	sd	s9,88(a0)
  2c:	07a53023          	sd	s10,96(a0)
  30:	07b53423          	sd	s11,104(a0)
  34:	0005b083          	ld	ra,0(a1)
  38:	0085b103          	ld	sp,8(a1)
  3c:	6980                	ld	s0,16(a1)
  3e:	6d84                	ld	s1,24(a1)
  40:	0205b903          	ld	s2,32(a1)
  44:	0285b983          	ld	s3,40(a1)
  48:	0305ba03          	ld	s4,48(a1)
  4c:	0385ba83          	ld	s5,56(a1)
  50:	0405bb03          	ld	s6,64(a1)
  54:	0485bb83          	ld	s7,72(a1)
  58:	0505bc03          	ld	s8,80(a1)
  5c:	0585bc83          	ld	s9,88(a1)
  60:	0605bd03          	ld	s10,96(a1)
  64:	0685bd83          	ld	s11,104(a1)
  68:	8082                	ret
exit code: 0
\end{lstlisting}

\section{Part 2: Run and Verify}
The complete output was checked against the expected sequence. After the three
start messages, the workers printed in the order C, A, B for every value from 0
through 99. Each worker exited after 100 increments. The transcript below is an
excerpt with 285 intermediate counter lines omitted. The full output is retained
in \texttt{evidence.txt}.

\begin{lstlisting}[caption={Actual xv6 run with the middle counter lines omitted}]
$ uthread
thread_a started
thread_b started
thread_c started
thread_c 0
thread_a 0
thread_b 0
thread_c 1
thread_a 1
thread_b 1
thread_c 2
thread_a 2
thread_b 2
[285 intermediate counter lines omitted]
thread_c 98
thread_a 98
thread_b 98
thread_c 99
thread_a 99
thread_b 99
thread_c: exit after 100
thread_a: exit after 100
thread_b: exit after 100
thread_schedule: no runnable threads
$ echo shell still runs
shell still runs
\end{lstlisting}

The final message follows the supplied skeleton. Once every worker has marked
itself free, there is no runnable slot, so the scheduler prints the message and
calls \texttt{exit(-1)}. The xv6 shell then accepts another command.

\begin{lstlisting}[caption={Checks on the full output, register layout, and source files}]
PASS: each worker printed 0 through 99 and exited after 100.
PASS: all 300 counter lines followed the c, a, b round-robin order.
PASS: all 14 register offsets match between C context and assembly.
PASS: kernel source files match the Lab 13 baseline byte for byte.
PASS: the xv6 shell accepted a command after uthread finished.
\end{lstlisting}

\section{Observe: Saved Registers and the Return Address}
\texttt{thread\_switch(old, new)} stores \texttt{ra}, \texttt{sp}, and
\texttt{s0} through \texttt{s11} in the old context, then loads those registers
from the new context. The stack pointer selects the other thread's stack.

A call to \texttt{thread\_switch} leaves a return address in \texttt{ra}. Saving
that register records where the suspended thread should resume. Restoring
another thread's \texttt{ra} and executing \texttt{ret} transfers control to its
saved continuation, just after its earlier switch call. A new thread has its
function address placed in \texttt{ra}, so the first return enters that function.

The \texttt{s} registers are callee-saved under the RISC-V calling convention.
The compiler handles caller-saved registers around the ordinary function call
boundary. The switch also preserves the return address and stack pointer needed
for the continuation. All three threads share the address space of one xv6
process and cooperate by yielding.

\section{Conclusion}
The context-switch routine saved and restored the register state required to
resume each worker on its own stack. The full output verified all 300 counter
lines and the three final values of 100. The implementation changes are in the
user program, its assembly routine, and the Makefile.

\appendix
\section{Source Code}

\begin{lstlisting}[caption={user/uthread.c},language=C]
#include "kernel/types.h"
#include "user/user.h"

#define FREE     0x0
#define RUNNING  0x1
#define RUNNABLE 0x2

#define STACK_SIZE 8192
#define MAX_THREAD 4

/* Layout matches the offsets used by uthread_switch.S. */
struct context {
  uint64 ra;
  uint64 sp;
  uint64 s0;
  uint64 s1;
  uint64 s2;
  uint64 s3;
  uint64 s4;
  uint64 s5;
  uint64 s6;
  uint64 s7;
  uint64 s8;
  uint64 s9;
  uint64 s10;
  uint64 s11;
};

struct thread {
  struct context ctx;
  char stack[STACK_SIZE] __attribute__((aligned(16)));
  int state;
};

struct thread all_thread[MAX_THREAD];
struct thread *current_thread;

extern void thread_switch(struct context *old, struct context *new);

void
thread_init(void)
{
  current_thread = &all_thread[0];
  current_thread->state = RUNNING;
}

void
thread_schedule(void)
{
  struct thread *t, *next_thread;

  next_thread = 0;
  t = current_thread + 1;
  for (int i = 0; i < MAX_THREAD; i++) {
    if (t >= all_thread + MAX_THREAD)
      t = all_thread;
    if (t->state == RUNNABLE) {
      next_thread = t;
      break;
    }
    t = t + 1;
  }

  if (next_thread == 0) {
    printf("thread_schedule: no runnable threads\n");
    exit(-1);
  }

  if (current_thread != next_thread) {
    struct thread *old = current_thread;
    next_thread->state = RUNNING;
    current_thread = next_thread;
    thread_switch(&old->ctx, &next_thread->ctx);
  }
}

void
thread_create(void (*func)(void))
{
  struct thread *t;

  for (t = all_thread; t < all_thread + MAX_THREAD; t++) {
    if (t->state == FREE)
      break;
  }
  if (t == all_thread + MAX_THREAD) {
    fprintf(2, "thread_create: no free thread slots\n");
    exit(1);
  }
  t->state = RUNNABLE;
  t->ctx.ra = (uint64)func;
  t->ctx.sp = (uint64)(t->stack + STACK_SIZE);
}

void
thread_yield(void)
{
  current_thread->state = RUNNABLE;
  thread_schedule();
}

volatile int a_started, b_started, c_started;
volatile int a_n, b_n, c_n;

void
thread_a(void)
{
  printf("thread_a started\n");
  a_started = 1;
  while (b_started == 0 || c_started == 0)
    thread_yield();

  for (int i = 0; i < 100; i++) {
    printf("thread_a %d\n", i);
    a_n += 1;
    thread_yield();
  }
  printf("thread_a: exit after %d\n", a_n);
  current_thread->state = FREE;
  thread_schedule();
}

void
thread_b(void)
{
  printf("thread_b started\n");
  b_started = 1;
  while (a_started == 0 || c_started == 0)
    thread_yield();

  for (int i = 0; i < 100; i++) {
    printf("thread_b %d\n", i);
    b_n += 1;
    thread_yield();
  }
  printf("thread_b: exit after %d\n", b_n);
  current_thread->state = FREE;
  thread_schedule();
}

void
thread_c(void)
{
  printf("thread_c started\n");
  c_started = 1;
  while (a_started == 0 || b_started == 0)
    thread_yield();

  for (int i = 0; i < 100; i++) {
    printf("thread_c %d\n", i);
    c_n += 1;
    thread_yield();
  }
  printf("thread_c: exit after %d\n", c_n);
  current_thread->state = FREE;
  thread_schedule();
}

int
main(void)
{
  a_started = b_started = c_started = 0;
  a_n = b_n = c_n = 0;
  thread_init();
  thread_create(thread_a);
  thread_create(thread_b);
  thread_create(thread_c);
  current_thread->state = FREE;
  thread_schedule();
  exit(0);
}
\end{lstlisting}

\begin{lstlisting}[caption={user/uthread\_switch.S}]
.text
.globl thread_switch
thread_switch:
  sd ra,   0(a0)
  sd sp,   8(a0)
  sd s0,  16(a0)
  sd s1,  24(a0)
  sd s2,  32(a0)
  sd s3,  40(a0)
  sd s4,  48(a0)
  sd s5,  56(a0)
  sd s6,  64(a0)
  sd s7,  72(a0)
  sd s8,  80(a0)
  sd s9,  88(a0)
  sd s10, 96(a0)
  sd s11,104(a0)

  ld ra,   0(a1)
  ld sp,   8(a1)
  ld s0,  16(a1)
  ld s1,  24(a1)
  ld s2,  32(a1)
  ld s3,  40(a1)
  ld s4,  48(a1)
  ld s5,  56(a1)
  ld s6,  64(a1)
  ld s7,  72(a1)
  ld s8,  80(a1)
  ld s9,  88(a1)
  ld s10, 96(a1)
  ld s11,104(a1)
  ret
\end{lstlisting}

\end{document}
